% !TEX root = ../main.tex

\chapter{系统实现}
\label{ch:implementation}

第三章详细阐述了模拟器在在线调度场景中的核心设计，以及其与离线模拟器在状态管理和预测演进过程中的区别。
本章进一步介绍全局调度器及性能预测器在大模型服务集群中的具体实现细节。

\section{全局调度器实现}
全局调度器采用Rust语言实现，利用其无垃圾回收带来的性能优势以及成熟的异步并发生态，以满足在线服务场景下调度关键路径对高效性和可靠性的要求。
系统基于Tokio运行时实现并发控制，从而能够对多个实例并发执行性能预测。性能预测组件被集成在全局调度器内部，并与vLLM推理引擎完成对接。
为保证不同策略之间比较的公平性，本文将调度决策过程与发送请求、接收响应等框架逻辑解耦，使不同策略在相同运行框架下仅需替换决策模块即可完成对比。
具体而言，基于性能预测的调度策略只需实现统一的Trait，并向框架暴露相应接口。
这种方式既实现了框架代码与策略逻辑的解耦，也为后续支持新的调度策略提供了便利。后续实验中使用的多种主流调度策略基线均基于这一框架实现。

\section{推理框架修改}
在现有集群框架上实现具备性能预测能力的调度策略时，第三章提出的“模拟器与推理实例保持状态同步”这一要求需要落到具体工程实现中，因此必须对推理框架代码进行一定修改。本文选用vLLM作为模型实例的推理框架。
% (#在这里介绍为什么需要低侵入性的修改框架代码，引出后面的修改内容)
如果对推理框架进行过于侵入式的修改，将使性能预测驱动的调度策略难以在现有集群服务中部署和维护。
因此，本文仅通过埋点记录实例调度器行为，并新增状态上报接口。相关改动主要集中在 Apiserver 和 Scheduler 两个组件中：
其中，Apiserver 暴露新的 SSE 接口，以维持长连接并持续上报状态信息；
对于 Scheduler，则在每轮构建批处理的 \texttt{Schedule} 函数中记录批处理行为及状态变化，例如KV Cache的分配与剔除，
并在批处理结束后、更新状态的 \texttt{update\_from\_output} 函数中记录生成 Token 数量以及请求完成情况，
随后统一交由 Apiserver 进行状态上报。上报信息的数据结构如（TODO!）所示。
在 vLLM 中，上述修改仅涉及少量新增逻辑，且基本不改变原有推理执行路径；对于其他推理框架，如 sgLang，也可通过类似方式完成集成。

\section{性能预测器}
除模拟器状态管理外，性能预测器是保证预测准确性的另一关键组件，同时也必须满足较低的在线预测时延要求。本文参考Vidur的设计思路，采用离线性能剖析与运行时查表相结合的方式进行性能预测。

具体来说，在离线性能剖析阶段，针对模型结构构建特定批次，将其送入模型执行前向传播，并记录各组成部分的执行时长。
在后续预测阶段，系统加载这些剖析结果，为每个模块构建“特征值到执行时长”的预测表，并通过随机森林对预测表中的缺失项进行填充，从而得到一张较为稠密的预测表。
运行时只需要加载各模块的预测表，根据模拟器构建出的批特征值查找对应执行时长，即可得到模型推理时间。这样便将性能剖析和高开销预测过程置于离线阶段完成，而在线预测只需执行若干次查表操作，可将预测时延降低到几次内存读取的量级，从而避免成为系统瓶颈。

在特征选定上，需要针对不同模型结构的执行行为采用不同特征来映射执行时长，同时尽量将每个组件使用的特征数量控制在一到两个，以避免特征过多导致预测表规模过大。
对于执行逻辑较简单的模块，如矩阵相乘（例如KV Projection、MLP等结构），执行时长主要与批次包含的Token数相关，因此建模时只需要使用Token数量这一特征。
但对于较为复杂的Attention层，前向推理的执行时长则更难刻画。尤其在PD共置部署下，Chunked Prefill和PD混合批等优化显著提高了批处理复杂度，因为推理时的动态性会构造出Prefill请求数、Decode请求数和各自KVCache大小均存在差异的批，给性能预测带来更大挑战。
针对混合批在Attention上的执行特性，本文将Prefill和Decode核函数分别建模以简化问题。前者使用KVCache大小和计算量（定义为应用Mask之后Attention算法所需的计算量，即Chunk\_size * kvcache\_size），后者使用Batch中的请求数和平均KVCache大小，两者分别通过特征映射得到执行时间。
这种建模方式既能够捕捉Prefill和Decode影响执行性能的关键因素，又能避免因特征过多而导致稠密预测表构建和动态查表过程难以实现。

由于加载的预测数据量较大，且对给定Batch进行性能预测的操作是无状态、只读的，因此性能预测模块只需维护一个预测器实例，即加载一次预测数据，
便能够在多个模拟器实例中并发进行性能预测。这既满足了在线调度场景对多实例并发预测的需求，也明显降低了内存开销。

\section{本章小结}
本章介绍了性能预测组件及全局调度框架的具体实现。
全局调度器使用Rust实现，将调度决策逻辑与框架逻辑解耦，通过统一的Trait接口支持多种调度策略的公平比较。
对vLLM推理框架的修改集中在Apiserver和Scheduler两个组件的埋点与状态上报，改动量小且易于迁移到其他推理框架。
性能预测器采用离线剖析与在线查表相结合的方式，针对不同模型组件选取关键特征构建预测表，将在线预测延迟控制在几次内存读取的量级，且预测器实例可在多个模拟器间共享以降低内存开销。
